% VAMOS Paper: Neurocomputing (Elsevier)
% =============================================================================
% Supplementary Material
% =============================================================================
\documentclass[final,5p,times,twocolumn]{elsarticle}

\usepackage{graphicx}
\graphicspath{{../generated/}}
\usepackage{booktabs}
\usepackage{amsmath}
\usepackage{amssymb}
\usepackage{multirow}
\usepackage{algorithm}
\usepackage{algpseudocode}
\usepackage{listings}
\usepackage{xcolor}
\usepackage{pifont}
\usepackage{enumitem}
\usepackage{xurl}
\newcommand{\cmark}{\ding{51}}
\newcommand{\pmark}{\ding{109}}
\newcommand{\xmark}{\ding{55}}
\usepackage[hidelinks]{hyperref}
\hypersetup{hypertexnames=false}
% Editorial changes relative to repository commit 2aa37a0.
% Keep revision markup separate from the scientific content and BibTeX records.
\usepackage{etoolbox}
\DeclareRobustCommand{\revision}[1]{\textcolor{red}{#1}}
\pdfstringdefDisableCommands{\def\revision#1{#1}\def\corref#1{}}

% Mark the two references requested for this revision without editing a .bbl.
\newcommand{\markrevisionreferences}{%
  \let\revisionoriginalbibitem\bibitem
  \renewcommand{\bibitem}[2][]{%
    \par
    \ifstrequal{##2}{HMM+22}{\color{red}}{%
      \ifstrequal{##2}{TCZ+17}{\color{red}}{\color{black}}}%
    \ifstrempty{##1}{\revisionoriginalbibitem{##2}}{%
      \revisionoriginalbibitem[##1]{##2}}}}

% The regenerated source table has revised repository links and source labels.
\makeatletter
\newenvironment{revisiontable}{%
  \begingroup
  \let\revisionoriginalfloatboxreset\@floatboxreset
  \def\@floatboxreset{\revisionoriginalfloatboxreset\color{red}}%
}{\endgroup}
\makeatother

\makeatletter
\g@addto@macro\UrlBreaks{\do\.\do\_\do\-\do\/}
\makeatother

% Code listing style
\lstset{
  language=Python,
  basicstyle=\ttfamily\scriptsize,
  keywordstyle=\bfseries\color{green!50!black},
  commentstyle=\itshape\color{teal},
  stringstyle=\color{red!70!black},
  identifierstyle=\color{black},
  showstringspaces=false,
  frame=lines,
  rulecolor=\color{gray!20},
  breaklines=true,
  numbers=left,
  numbersep=5pt,
  numberstyle=\tiny\color{black},
  captionpos=b,
  xleftmargin=.5em,
  framexleftmargin=.5em,
  aboveskip=.5em,
  belowskip=.5em
}

\renewcommand{\lstlistingname}{Code}% Listing -> Code

\newcommand{\VAMOS}{VAMOS}
\newcommand{\HV}{\mathrm{HV}}
\newcommand{\Intdomain}[1]{[#1]_{\mathbb{Z}}}
\newcommand{\Realdomain}[1]{[#1]_{\mathbb{R}}}

\journal{\revision{Neurocomputing}}
\biboptions{sort&compress}

% Independent numbering distinguishes this document from the main manuscript.
\renewcommand{\thesection}{\revision{S\arabic{section}}}
\renewcommand{\thetable}{\revision{S\arabic{table}}}
\renewcommand{\thefigure}{\revision{S\arabic{figure}}}
\renewcommand{\theequation}{\revision{S\arabic{equation}}}

\begin{document}

\begin{frontmatter}

\title{VAMOS: A Framework for Fast, Configurable and Accessible Multiobjective Optimization in Python\\Supplementary Material}

\author[urjc]{Nicol\'{a}s R. Uribe\corref{cor1}}
\ead{nicolas.rodriguez@urjc.es}
\author[urjc]{Alberto Herr\'{a}n}
\ead{alberto.herran@urjc.es}
\author[uma]{Antonio J. Nebro}
\ead{ajnebro@uma.es}
\author[tecnalia,upv]{Javier Del Ser}
\ead{javier.delser@tecnalia.com}
\author[urjc]{J. Manuel Colmenar}
\ead{josemanuel.colmenar@urjc.es}

\cortext[cor1]{Corresponding author}
\address[urjc]{Department of Computer Sciences, Universidad Rey Juan Carlos, M\'{o}stoles, 28933 Madrid, Spain}
\address[uma]{Department of Lenguajes y Ciencias de la Computaci\'{o}n, ITIS Software, University of M\'{a}laga, 29071 M\'{a}laga, Spain}
\address[tecnalia]{TECNALIA, Basque Research \& Technology Alliance (BRTA), Derio, 48160, Spain}
\address[upv]{University of the Basque Country (UPV/EHU), Leioa, 48940, Spain}

\end{frontmatter}

This Supplementary Material accompanies the \revision{Neurocomputing} manuscript and contains the extended convergence, solution-quality, runtime, statistical, configurability, and accessibility results referenced there. \revision{The main manuscript presents the architecture, experimental protocol, primary runtime comparisons, and conclusions. Sections, tables, figures, and equations in this document use an S prefix and are numbered independently of the main manuscript.}

%==============================================================================
\section{Convergence Analysis}
\label{app:convergence}
%==============================================================================

\begin{figure*}[!t]
  \centering
  \includegraphics[width=\textwidth]{../generated/figures/convergence.png}
  \caption{Convergence of normalized hypervolume over 50{,}000 evaluations for \VAMOS{} (Numba), \textit{pymoo}, and \textit{jMetalPy} on DTLZ2 (3 objectives, 12 variables), WFG4 (2 objectives, 24 variables), and ZDT1 (2 objectives, 30 variables). Solid/dashed lines show median hypervolume values across 30 seeds; shaded regions indicate the interquartile range. All three frameworks follow nearly identical convergence trajectories, confirming that the semantic alignment protocol produces equivalent algorithmic behavior.}
  \label{fig:convergence}
\end{figure*}

To verify that runtime speedups do not come at the expense of convergence rate, we track normalized hypervolume at 1{,}000-evaluation intervals over the full 50{,}000-evaluation budget. Fig.~\ref{fig:convergence} compares \VAMOS{}, \textit{pymoo}\revision{,} and \textit{jMetalPy} on three representative problems: DTLZ2 (3 objectives, 12 variables), WFG4 (2 objectives, 24 variables), and ZDT1 (2 objectives, 30 variables). The convergence curves are nearly indistinguishable across all three frameworks, with overlapping interquartile ranges across all 30 seeds. All frameworks reach the same final quality at the same evaluation counts, confirming that the performance gains observed in runtime measurements are not accompanied by degraded convergence behavior.


%==============================================================================
\section{Solution Quality Analysis}
\label{app:quality}
%==============================================================================
The runtime advantages reported \revision{in the main manuscript} are only meaningful if \VAMOS{} produces solutions of comparable quality. To verify this, we compare normalized hypervolume (HV) across frameworks using the protocol described in the \revision{Performance} section of the main \revision{manuscript:} for each problem we compute the median HV across 30 seeds, and then summarize these per-problem medians by family using the median and interquartile range (IQR). Table~\ref{tab:frameworks_hv} presents the results for the generational Non-dominated Sorting Genetic Algorithm II (NSGA-II).

\begin{table*}[htbp]
\centering
%\scriptsize
%\setlength{\tabcolsep}{3pt}
\caption{Normalized hypervolume summary (median (IQR)) by problem family across frameworks}
\label{tab:frameworks_hv}
\begin{tabular}{@{}lccccc@{}}
\toprule
\textbf{Framework} & \textbf{ZDT} & \textbf{DTLZ} & \textbf{WFG} & \textbf{ZCAT} & \textbf{Overall} \\
\midrule
VAMOS & \textbf{0.991 (0.013)} & 0.836 (0.121) & \textbf{0.934 (0.138)} & --- & \textbf{0.934 (0.157)} \\
\textit{pymoo} & \textbf{0.991 (0.010)} & 0.829 (0.148) & 0.846 (0.220) & --- & 0.921 (0.156) \\
\textit{jMetalPy} & \textbf{0.991 (0.014)} & 0.835 (0.110) & 0.845 (0.223) & --- & 0.924 (0.159) \\
\textit{DEAP} & \textbf{0.991 (0.013)} & \textbf{0.839 (0.138)} & 0.839 (0.224) & --- & 0.924 (0.156) \\
\textit{Platypus} & \textbf{0.991 (0.010)} & 0.824 (0.177) & 0.919 (0.137) & --- & 0.919 (0.160) \\
\bottomrule
\end{tabular}
\end{table*}

Across ZDT, all frameworks attain nearly identical normalized HV, confirming that benchmark definitions and operator semantics are effectively aligned. For DTLZ, medians are tightly clustered across frameworks. WFG exhibits the largest dispersion; notably, \VAMOS{} achieves the highest median on WFG (0.934 vs.\ 0.846 for \textit{pymoo}\revision{),} which may reflect subtle differences in WFG parameterization or boundary handling across implementations. Together with the runtime results, these summaries confirm that \VAMOS{} achieves comparable or superior solution quality while delivering faster runtimes under the fixed-reference-front protocol.

Table~\ref{tab:frameworks_hv_variants} extends this analysis to the NSGA-II steady-state and external-archive variants. For the steady-state variant, the overall normalized HV is effectively identical across frameworks (all three reach 0.989 overall), although \VAMOS{} attains a clearly higher median on WFG (0.939 vs.\ 0.849 for \textit{pymoo} and 0.848 for jMetalPy). For the archive variant, the three frameworks remain closely matched across all families: ZDT and ZCAT are virtually identical, DTLZ differs by at most 0.003, and the WFG medians lie in the narrow 0.846--0.852 range.

\begin{table*}[htbp]
\centering
%\scriptsize
%\setlength{\tabcolsep}{3pt}
\caption{Normalized hypervolume summary (median (IQR)) by problem family for NSGA-II variants}
\label{tab:frameworks_hv_variants}
\begin{tabular}{@{}llccccc@{}}
\toprule
\textbf{Algorithm} & \textbf{Framework} & \textbf{ZDT} & \textbf{DTLZ} & \textbf{WFG} & \textbf{ZCAT} & \textbf{Overall} \\
\midrule
NSGA-II (SS) & VAMOS & 0.993 (0.009) & 0.846 (0.091) & 0.939 (0.138) & 0.994 (0.004) & 0.989 (0.055) \\
 & \textit{pymoo} & 0.993 (0.008) & 0.847 (0.160) & 0.849 (0.218) & 0.994 (0.004) & 0.989 (0.073) \\
 & \textit{jMetalPy} & 0.993 (0.009) & 0.849 (0.101) & 0.848 (0.153) & 0.994 (0.004) & 0.989 (0.073) \\
\midrule
NSGA-II (Archive) & VAMOS & 0.993 (0.010) & 0.855 (0.106) & 0.846 (0.219) & 0.993 (0.005) & 0.987 (0.079) \\
 & \textit{pymoo} & 0.993 (0.009) & 0.853 (0.139) & 0.852 (0.223) & 0.993 (0.005) & 0.988 (0.080) \\
 & \textit{jMetalPy} & 0.993 (0.010) & 0.856 (0.104) & 0.846 (0.149) & 0.993 (0.005) & 0.987 (0.073) \\
\bottomrule
\end{tabular}
\end{table*}

%\input{../generated/tables/frameworks_igd_merged.tex}

%==============================================================================
\section{Statistical Analysis}
\label{app:stats_main}
%==============================================================================

The descriptive summaries in the previous section suggest quality equivalence, but median and IQR alone cannot determine whether observed differences are statistically significant. We therefore apply Wilcoxon signed-rank tests (paired by seed) comparing \VAMOS{} against each competitor on every problem, with Holm step-down correction for familywise error control ($\alpha=0.05$). As a measure of practical significance, we report the Vargha--Delaney $\hat{A}_{12}$ effect size~\cite{VD00} alongside each test ($\hat{A}_{12}>0.5$ favors \VAMOS{}; thresholds: negligible${}<0.56$, small${}<0.64$, medium${}<0.71$, large${\ge}0.71$). We also compute paired-bootstrap 90\% confidence intervals (CIs) for the relative normalized-HV difference $\Delta = (\HV_{\text{VAMOS}} - \HV_{\text{fw}})/\HV_{\text{fw}}$, and declare \emph{equivalence} when the entire CI falls within $\pm 1\%$.

Table~\ref{tab:stats_hypervolume_nsgaii} presents the per-problem Wilcoxon results for the NSGA-II baseline. Of 41 problems, only one (ZDT6) shows a statistically significant difference in normalized HV after Holm correction, and the magnitude is small ($\Delta \approx 0.4\%$). Table~\ref{tab:hv_equivalence_summary_nsgaii} summarizes the equivalence analysis: \VAMOS{} is equivalent to \textit{pymoo} on 37/41 problems and non-inferior on 38/41, with a median $\Delta$HV of $-0.01\%$. Similar patterns hold for \textit{jMetalPy} and the other frameworks. Table~\ref{tab:hv_robustness_summary_nsgaii} details robustness, and Fig.~\ref{fig:forest_ci_nsgaii} visualizes per-problem CIs as a forest plot. Detailed NSGA-II statistical tables are provided in Section~\ref{app:stats_nsgaii}, and corresponding analyses for SMS-EMOA and MOEA/D are provided in Section~\ref{app:stats}.

%==============================================================================
\section{Additional \revision{Runtime Summaries}}
\label{app:variants}
%==============================================================================
\revision{The main manuscript} reports \revision{the family-level median runtimes} for \revision{the} NSGA-II \revision{steady-state} and \revision{external-archive variants, SMS-EMOA, and MOEA/D in its Cross-Framework Comparison subsection. Table~\ref{tab:runtime_family_averages} complements those results with their unweighted arithmetic mean across} the \revision{four benchmark families (ZDT, DTLZ, WFG,} and \revision{ZCAT). Each family contributes equally; these values are not medians pooled across all problems or seeds.}

\begin{table}[htbp]
\color{red}
\centering
\caption{Unweighted mean of the four family-level median runtimes (seconds).}
\label{tab:runtime_family_averages}
\begin{tabular}{lrrr}
\toprule
\textbf{Algorithm} & \textbf{VAMOS} & \textbf{pymoo} & \textbf{jMetalPy} \\
\midrule
NSGA-II (SS) & \textbf{477.40} & 5989.77 & 2285.27 \\
NSGA-II (Archive) & \textbf{56.91} & 2914.52 & 1769.51 \\
SMS-EMOA & \textbf{391.94} & 2618.83 & 1184.56 \\
MOEA/D & \textbf{295.98} & 2057.00 & 1338.74 \\
\bottomrule
\end{tabular}
\end{table}

\subsection{Memory footprint}

Peak resident-set-size (RSS) was measured during NSGA-II runs on three representative problems (ZDT1, DTLZ2, WFG4) using 50{,}000 evaluations and 5 seeds per framework. \VAMOS{} achieves a median peak RSS of 0.28 megabytes (MB), compared to 0.84\,MB for \textit{pymoo} ($3\times$ higher) and 0.50\,MB for \textit{jMetalPy} ($1.8\times$ higher). These measurements are consistent with the array-native representation reducing per-individual object overhead.

\begin{figure}[htbp]
\centering
\includegraphics[width=\columnwidth]{../generated/figures/memory_comparison.png}
\caption{Peak memory footprint (MB) during NSGA-II optimization on three representative problems (50{,}000 evaluations, median of 5 seeds). \VAMOS{} consistently uses less memory than both \textit{pymoo} and \textit{jMetalPy} across the measured problems.}
\label{fig:memory_comparison}
\end{figure}

\section{Detailed statistical analysis for NSGA-II}
\label{app:stats_nsgaii}
This section reports per-problem Wilcoxon signed-rank tests, bootstrap equivalence analyses, robustness summaries, and bootstrap confidence intervals for the NSGA-II baseline discussed in \revision{Section~\ref{app:stats_main}} of \revision{this Supplementary Material.}
\begingroup
\setlength{\tabcolsep}{2pt}
\input{../generated/tables/stats_nsgaii.tex}
\endgroup

\section{Reference fronts and detailed benchmark results}
\label{app:detailed}
%==============================================================================

The \revision{official} repository \revision{at \url{https://github.com/vamos-optimization/VAMOS}} includes \revision{the benchmark resources and retained research inputs. Dense} reference Pareto fronts \revision{are used} for all benchmark problems \revision{for} normalized hypervolume computation. These fronts are generated analytically when closed forms are available (ZDT and most DTLZ problems) and by dense sampling followed by non-dominated filtering for cases where the Pareto front is disconnected or defined implicitly (DTLZ7 and WFG). For WFG2 we use a higher sampling density to stabilize HV normalization.

Fig.~\ref{fig:runtime_heatmap} presents per-problem median runtimes as a dual-panel heatmap. The left panel compares VAMOS backends (Numba, MooCore, NumPy); the right panel compares VAMOS against \textit{pymoo} and jMetalPy. Color intensity indicates runtime on a logarithmic scale, with annotated values in each cell and bold values marking the fastest per problem.

\begin{figure*}[htbp]
  \centering
  \includegraphics[width=0.95\textwidth]{../generated/figures/runtime_heatmap.png}
  \caption{Per-problem median runtime (seconds) for VAMOS backends (left) and cross-framework comparison (right). Color scale is logarithmic; bold values indicate the fastest in each row. Problems are grouped by family (ZDT, DTLZ, WFG).}
  \label{fig:runtime_heatmap}
\end{figure*}

%==============================================================================
\section{Statistical analysis for SMS-EMOA and MOEA/D}
\label{app:stats}
%==============================================================================

This section reports Wilcoxon signed-rank tests, bootstrap equivalence analyses, and robustness summaries for SMS-EMOA and MOEA/D, following the same methodology as \revision{Section~\ref{app:stats_main}} of \revision{this Supplementary Material.}

\subsection{SMS-EMOA}

Tables~\ref{tab:stats_hypervolume_smsemoa}--\ref{tab:hv_robustness_summary_smsemoa} present per-problem Wilcoxon tests, equivalence summaries, and robustness summaries for SMS-EMOA. Fig.~\ref{fig:forest_ci_smsemoa} visualizes per-problem bootstrap confidence intervals.

\begingroup
\setlength{\tabcolsep}{2pt}
\input{../generated/tables/stats_smsemoa.tex}
\endgroup

\subsection{MOEA/D}

Tables~\ref{tab:stats_hypervolume_moead}--\ref{tab:hv_robustness_summary_moead} and Fig.~\ref{fig:forest_ci_moead} present the corresponding analyses for MOEA/D.

\begingroup
\setlength{\tabcolsep}{2pt}
\input{../generated/tables/stats_moead.tex}
\endgroup

%==============================================================================
\section{Accessibility Proxy Protocol}
\label{app:accessibility_protocol}
%==============================================================================

To complement the accessibility discussion in the main paper, we evaluate three onboarding tasks: (i)~run NSGA-II on a built-in benchmark, (ii)~define and optimize a small custom two-objective problem, and (iii)~move from defaults to explicit operator-level configuration. For each framework, we selected a canonical snippet from the shortest documented path available in March~2026. For \VAMOS{}, the snippets \revision{illustrate the same tasks as} the public API examples \revision{in} the manuscript; for \textit{pymoo} and jMetalPy, the canonical sources are the public documentation pages listed in Table~\ref{tab:accessibility_proxy_details}.

Metrics count only non-empty, non-comment lines and explicit \texttt{import}/\texttt{from} statements. Plotting, result serialization, and explanatory comments are excluded because they do not affect the onboarding burden of the optimization call itself. The underlying artifact file used to generate the accessibility-proxy table in the main paper and Table~\ref{tab:accessibility_proxy_details} below is \texttt{paper/accessibility\_proxy\_snippets.json}.

\begin{revisiontable}
\input{../generated/tables/accessibility_proxy_details.tex}
\end{revisiontable}


%==============================================================================
\section{Code \revision{and Data Availability}}
\label{app:code}
%==============================================================================

\revision{The official repository is \url{https://github.com/vamos-optimization/VAMOS}. It provides source code, examples, tuning workflows,} and \revision{the retained benchmark inputs. The publication workflow is documented in \texttt{paper/README.md}, and the retained research inputs are described in \texttt{experiments/REFERENCE\_RESULTS.md}. The experiments reported in the main manuscript used the internal development build identified in its Runtime Measurement subsection; the public 1.0.0 release is a subsequent software release.}

\revision{The data inventory explains the measurement scope of each experiment. In \texttt{scaling\_vectorization.csv}, runtime is recorded for all 220 observations. Hypervolume calculation is explicitly disabled for the 40 DTLZ2 objective-count observations (2, 3, 5, and 8 objectives; NumPy and Numba; five seeds per combination). The JIT timing-policy field is inapplicable to the 95 NumPy observations. These empty cells follow the experiment's recording conventions. The steady-state CSV also contains one additional jMetalPy observation on ZCAT1 with seed 0 and 200 evaluations. The reviewer summary keeps this observation separate from the 50,000-evaluation campaign.}


%==============================================================================
\section{Parameter Spaces for pymoo and jMetalPy}
\label{app:parameter_spaces}
%==============================================================================

Tables~\ref{tab:parameter_space_pymoo} and~\ref{tab:parameter_space_jmetalpy} detail the configurable parameter spaces for NSGA-II in \textit{pymoo} and \textit{jMetalPy}, respectively.

\begin{table*}[htbp]
\centering
\caption{Parameter space for NSGA-II in \textit{pymoo} (real-valued encoding). Parameters in italics are conditional, active only when their corresponding operator is selected. Values in \textbf{bold} indicate defaults. The total number of configurable parameters is 20 (11 always active, 9~conditional). Acronyms in the table denote Latin hypercube sampling (LHS), parent-centric crossover (PCX), and differential evolution crossover (DEX).}
\label{tab:parameter_space_pymoo}
\resizebox{\textwidth}{!}{%
\begin{tabular}{lll}
\toprule
\textbf{Component} & \textbf{Value / Strategy} & \textbf{Conditional Parameters} \\ \midrule
\multirow{2}{*}{General} & Pop.\ Size: integer; default $\mathbf{100}$ & -- \\
 & Offspring Count (\texttt{n\_offsprings}): integer; default = Pop.\ Size & -- \\ \midrule
\multirow{1}{*}{Initialization} & \{\textbf{Random}, LHS\} & -- \\ \midrule
\multirow{3}{*}{Selection} & \{\textbf{Tournament}, Random\} & -- \\
 & \textit{Tournament Size} (\texttt{pressure}): integer; default $\mathbf{2}$ & \textit{Selection} = Tournament \\
 & Crowding Metric $\in$ \{\textbf{cd}, pcd, ce, mnn, 2nn\} & -- \\ \midrule
\multirow{4}{*}{Crossover} & \multicolumn{2}{l}{\textit{Global:} Probability $\in \Realdomain{0, 1}$; default $\mathbf{0.9}$} \\
 & \textbf{SBX} & \textit{Dist.\ Index} $\eta \in \Realdomain{3, 30}$; default $\mathbf{15}$. \textit{Per-var.\ prob.}\ $\in \Realdomain{0.1, 0.9}$; default $\mathbf{0.5}$ \\
 & PCX & $\eta \in \Realdomain{0.01, 0.3}$, $\zeta \in \Realdomain{0.01, 0.3}$ \\
 & DEX & $CR \in \Realdomain{0, 1}$; default $0.7$. $F$: real or \texttt{None} \\ \midrule
\multirow{3}{*}{Mutation} & \multicolumn{2}{l}{\textit{Global:} Prob.\ $\in \Realdomain{0, 1}$; default $\mathbf{0.9}$. Per-var.\ prob.; default $\min(0.5, 1/n)$} \\
 & \textbf{Polynomial} & \textit{Dist.\ Index} $\eta \in \Realdomain{3, 30}$; default $\mathbf{20}$ \\
 & Gaussian & $\sigma \in \Realdomain{0.01, 0.25}$; default $0.1$ \\ \midrule
\multirow{1}{*}{Repair} & \{\textbf{clip}, bounce-back, to-bound\} & -- \\ \bottomrule
\end{tabular}%
}
\end{table*}

\begin{table*}[htbp]
\centering
\caption{Parameter space for NSGA-II in \textit{jMetalPy} (real-valued encoding). Parameters in italics are conditional, active only when their corresponding operator is selected. Values in \textbf{bold} indicate defaults. The total number of configurable parameters is 29 (10 always active, 19~conditional). Acronyms in the table denote unimodal normal distribution crossover (UNDX) and differential evolution (DE).}
\label{tab:parameter_space_jmetalpy}
%\resizebox{\textwidth}{!}{%
\begin{tabular}{lll}
\toprule
\textbf{Component} & \textbf{Value / Strategy} & \textbf{Conditional Parameters} \\ \midrule
\multirow{5}{*}{General} & Pop.\ Size: integer & -- \\
 & Offspring Pop.\ Size: integer; $= 1$ for steady-state & -- \\
 & Use External Archive $\in$ \{True, \textbf{False}\} & -- \\
 & \textit{Archive Type} $\in$ \{NonDominated, CrowdingDist., DistanceBased\} & \textit{Archive} = True \\
 & \textit{Archive Max Size}: integer & \textit{Archive} = True, bounded \\ \midrule
\multirow{1}{*}{Initialization} & \{\textbf{Random}, Injector\} & -- \\ \midrule
\multirow{2}{*}{Selection} & \{\textbf{Binary Tournament}, Tournament, Random\} & -- \\
 & \textit{Tournament Size}: integer $\geq 2$ & \textit{Selection} = Tournament \\ \midrule
\multirow{7}{*}{Crossover} & \multicolumn{2}{l}{\textit{Global:} Probability $\in \Realdomain{0, 1}$} \\
 & \textbf{SBX} & \textit{Dist.\ Index}; default $\mathbf{20}$ \\
 & BLX-$\alpha$ & $\alpha$; default $0.5$ \\
 & BLX-$\alpha\beta$ & $\alpha$, $\beta$; default $0.5$ each \\
 & Arithmetic & -- \\
 & UNDX & $\zeta$; default $0.5$. $\eta$; default $0.35$ \\
 & DE & $CR \in \Realdomain{0, 1}$, $F \in \Realdomain{0, 2}$, $K \in \Realdomain{0, 1}$ \\ \midrule
\multirow{7}{*}{Mutation} & \multicolumn{2}{l}{\textit{Global:} Probability $\in \Realdomain{0, 1}$; conventionally $1/n$} \\
 & \textbf{Polynomial} & \textit{Dist.\ Index}; default $\mathbf{20}$ \\
 & Simple Random & -- \\
 & Uniform & \textit{Perturbation}; default $0.5$ \\
 & Non-Uniform & \textit{Perturbation}, \textit{Max Iterations} \\
 & L\'evy Flight & $\beta$; default $1.5$. \textit{Step size}; default $0.01$ \\
 & Power Law & $\delta$; default $1.0$ \\ \midrule
\multirow{1}{*}{Repair} & \{\textbf{Clamp}, Random Uniform, Reflective, Bound Swap\} & -- \\ \bottomrule
\end{tabular}%
%}
\end{table*}

%==============================================================================
% Bibliography
%==============================================================================
\bibliographystyle{elsarticle-num}
\bibliography{references}

\end{document}
